\documentclass[10pt, letterpaper]{article}

% ── Paquetes ──────────────────────────────────────────────────────────────────
\usepackage[utf8]{inputenc}
\usepackage[T1]{fontenc}
\usepackage{babel}
\usepackage[top=1.1cm, bottom=1.0cm, left=1.3cm, right=1.3cm]{geometry}
\usepackage{lmodern}
\usepackage{microtype}
\usepackage{parskip}
\usepackage{enumitem}
\usepackage{titlesec}
\usepackage{hyperref}
\usepackage{xcolor}
\usepackage{tabularx}
\usepackage{array}
\usepackage{booktabs}
\usepackage{multicol}
\usepackage{setspace}

% ── Colores ───────────────────────────────────────────────────────────────────
\definecolor{azuloscuro}{RGB}{10, 40, 80}
\definecolor{azulmedio}{RGB}{30, 80, 150}
\definecolor{gristexto}{RGB}{55, 55, 55}
\definecolor{grisclaro}{RGB}{235, 238, 242}
\definecolor{lineagris}{RGB}{180, 190, 205}

% ── Hiperlinks ────────────────────────────────────────────────────────────────
\hypersetup{colorlinks=true, urlcolor=azulmedio, linkcolor=azulmedio}

% ── Tipografía y espaciado ────────────────────────────────────────────────────
\setlength{\parskip}{0pt}
\setlength{\parindent}{0pt}
\renewcommand{\baselinestretch}{1.0}

% ── Sección con línea ─────────────────────────────────────────────────────────
\titleformat{\section}
  {\color{azuloscuro}\normalfont\bfseries\normalsize}
  {}{0em}{\MakeUppercase}[{\color{lineagris}\titlerule[0.5pt]}]
\titlespacing{\section}{0pt}{5pt}{3pt}

% ── Lista compacta ────────────────────────────────────────────────────────────
\setlist[itemize]{
  leftmargin=1.2em,
  itemsep=0.3pt,
  topsep=1pt,
  parsep=0pt,
  label={\color{azulmedio}\textbullet}
}

% ── Comandos de CV ────────────────────────────────────────────────────────────
\newcommand{\cventry}[4]{%
  \vspace{2pt}%
  {\color{azuloscuro}\textbf{#1}}\hfill{\small\color{gristexto}\textit{#2}}\\[-3pt]
  {\small\color{azulmedio}\textit{#3}}\hfill{\small\color{gristexto}#4}%
  \vspace{1pt}%
}

\newcommand{\cvedu}[4]{%
  \vspace{1pt}%
  {\color{azuloscuro}\textbf{#1}}\hfill{\small\color{gristexto}\textit{#2}}\\[-3pt]
  {\small\color{gristexto}#3}\hfill{\small\color{azulmedio}\textit{#4}}%
  \vspace{1pt}%
}

\newcommand{\skillcat}[2]{%
  \textbf{\color{azuloscuro}#1:} {\small\color{gristexto}#2}\\[0.8pt]%
}

% ═════════════════════════════════════════════════════════════════════════════
\begin{document}
\pagestyle{empty}

% ── ENCABEZADO ────────────────────────────────────────────────────────────────
\begin{center}
  {\LARGE\bfseries\color{azuloscuro} JAIR MOLINA ARCE}\\[3pt]
  {\small\color{azulmedio}\textbf{Ingeniero en Automatización Industrial · Sistemas Embebidos · Investigador Aeroespacial}}\\[4pt]
  {\footnotesize\color{gristexto}
    \href{mailto:ingjairmolina@gmail.com}{ingjairmolina@gmail.com}\enspace$\cdot$\enspace
    +52 56 52 64 61 08\enspace$\cdot$\enspace
    Ciudad de México, México\enspace$\cdot$\enspace
    ORCID: \href{https://orcid.org/0009-0009-6732-8100}{0009-0009-6732-8100}\enspace$\cdot$\enspace
    CVU CONACYT: 1340773
  }\\[1pt]
  {\footnotesize\color{azulmedio}
    \href{https://codeaerospace.com}{codeaerospace.com}\enspace$\cdot$\enspace
    \href{https://codeaerospace.space}{codeaerospace.space}\enspace$\cdot$\enspace
    \href{https://codeaerospace.tech}{codeaerospace.tech}
  }
\end{center}

\vspace{-2pt}
{\color{lineagris}\rule{\linewidth}{0.5pt}}
\vspace{2pt}

% ── PERFIL PROFESIONAL ────────────────────────────────────────────────────────
\section{Perfil Profesional}

{\small\color{gristexto}
Ingeniero Mecánico con Maestría en curso (IPN) en \textbf{Control de Sistemas Dinámicos e Instrumentación}. Co-fundador de \textbf{Co.De Aerospace}: organización de investigación y desarrollo que opera una estación terrena SDR para seguimiento de satélites (\textit{Orbit Eye}), una plataforma de investigación 6G con IA y asesoría para el primer hábitat análogo de México. Experiencia en \textbf{DAQ} industrial, controladores \textbf{PID} digitales, \textbf{MATLAB/Simulink} y lenguajes \textbf{IEC~61131-3}. Docente universitario en ingeniería térmica y sistemas aeroespaciales. Publicaciones científicas internacionales. Inglés avanzado. \textbf{Disponibilidad inmediata para radicarse en EE.UU.}
}

% ── EXPERIENCIA ───────────────────────────────────────────────────────────────
\section{Experiencia Profesional}

% ── Co.De Aerospace
\cventry{Co.De Aerospace — Co-fundador \& Ingeniero Líder de Tecnología}{2023 – presente}{\href{https://codeaerospace.com}{codeaerospace.com} · \href{https://codeaerospace.space}{codeaerospace.space} · \href{https://codeaerospace.tech}{codeaerospace.tech}}{Remoto / México}
\begin{itemize}
  \item Desarrollé \textbf{Co.De Orbit Eye}: estación terrena SDR para seguimiento en tiempo real de satélites en órbita baja (LEO), con interfaz web de monitoreo y visualización orbital (\href{https://codeaerospace.space}{codeaerospace.space}).
  \item Diseñé \textbf{ResearchAI}: plataforma de investigación en telecomunicaciones \textbf{6G} potenciada con IA, orientada a comunicaciones satelitales de próxima generación (\href{https://codeaerospace.tech}{codeaerospace.tech}).
  \item Lideré colaboraciones técnicas en \textbf{ciberseguridad satelital}: protocolos de comunicación segura y hardening de sistemas de control en tierra.
  \item Brindé asesoría técnica para el diseño e implementación del \textbf{primer hábitat análogo de México}, incluyendo sistemas de instrumentación y control ambiental.
  \item Coordiné demostraciones tecnológicas y divulgación científica ante medios nacionales (TV Azteca, 2024) e instancias internacionales.
\end{itemize}

\vspace{2pt}

% ── Banco de Pruebas IPN
\cventry{Sistemas Embebidos e Instrumentación — Banco de Pruebas de Propulsión}{2022 – 2024}{Instituto Politécnico Nacional (IPN)}{Ciudad de México}
\begin{itemize}
  \item Diseñé e implementé sistema \textbf{DAQ} multicanal con sensores industriales de presión, temperatura y flujo para banco de pruebas aeroespacial.
  \item Programé microcontroladores \textbf{ESP32} y \textbf{STM32} en \textbf{C/C++} para adquisición y preprocesamiento de señales en tiempo real.
  \item Implementé controladores \textbf{PID digitales} y modelos de sistemas dinámicos en \textbf{MATLAB/Simulink}, logrando estabilidad en lazo cerrado.
  \item Diseñé protocolo de transmisión \textbf{MQTT/WiFi} para monitoreo remoto continuo; procesé datos con \textbf{Python} y \textbf{LabVIEW}.
\end{itemize}

\vspace{2pt}

% ── JobHunter
\cventry{Plataforma de Automatización con IA: \textit{JobHunter}}{2024 – 2025}{Proyecto Independiente}{Remoto}
\begin{itemize}
  \item Backend \textbf{FastAPI + Python} con integración Groq AI (LLM); frontend \textbf{Next.js/TypeScript} con dashboard en tiempo real.
  \item Bot Telegram para notificaciones y control de flujos automatizados, reduciendo intervención manual en +80\%.
\end{itemize}

\vspace{2pt}

% ── IoT
\cventry{Plataforma Web Full-Stack para Monitoreo IoT}{2023 – 2024}{Proyecto Independiente / IPN}{Ciudad de México}
\begin{itemize}
  \item Plataforma de monitoreo remoto de sensores con \textbf{Flask + SQLAlchemy}, APIs RESTful, \textbf{Docker}, Prometheus, MySQL y Redis.
\end{itemize}

\vspace{2pt}

% ── Docente UNII
\cventry{Docente — Ingeniería Térmica y Modelado de Sistemas Aeroespaciales}{2025 – presente}{UNII — Universidad de la Innovación, Aguascalientes}{Remoto}
\begin{itemize}
  \item Imparto cursos de ingeniería térmica aplicada y modelado/simulación de sistemas aeroespaciales a nivel licenciatura en modalidad remota.
\end{itemize}

\vspace{2pt}

% ── Profuturo
\cventry{Capacitador — Excel Avanzado para Análisis de Datos}{Marzo 2026}{Profuturo}{Ciudad de México}
\begin{itemize}
  \item Diseñé e impartí taller intensivo de Excel avanzado (tablas dinámicas, Power Query, dashboards, automatización con macros VBA).
\end{itemize}

\vspace{2pt}

% ── ENBA
\cventry{Docente — Tecnologías Disruptivas (Posgrado)}{2024}{ENBA – IPN}{Ciudad de México}
\begin{itemize}
  \item Impartí curso de VR/AR aplicado a capacitación técnica industrial a nivel posgrado.
\end{itemize}

% ── FORMACIÓN ACADÉMICA ───────────────────────────────────────────────────────
\section{Formación Académica}

\cvedu{Maestría en Tecnologías Avanzadas}{2024 – presente}{IPN — Línea: Control de sistemas dinámicos e instrumentación}{En curso}
\vspace{2pt}
\cvedu{Ingeniería Mecánica}{2017 – 2023}{IPN — Térmica, instrumentación, sistemas embebidos, bancos de prueba}{Titulado}
\vspace{2pt}
\cvedu{Bachillerato Técnico en Informática}{2014 – 2017}{Colegio de Bachilleres Plantel 1 ``El Rosario''}{Concluido}

% ── COMPETENCIAS TÉCNICAS ─────────────────────────────────────────────────────
\section{Competencias Técnicas}

{\small
\skillcat{Automatización / Control}{\textbf{Ladder Logic, Structured Text, Function Block Diagram} (IEC~61131-3), \textbf{PID} digital, modelado sistemas dinámicos, \textbf{MATLAB/Simulink}}
\skillcat{Embebidos / IoT}{\textbf{ESP32}, \textbf{STM32}, Arduino, \textbf{C/C++}, MQTT, HTTP/REST, WiFi, UART, SPI, I2C}
\skillcat{Instrumentación / DAQ}{Sensores presión/temperatura/flujo, acondicionamiento señales, \textbf{DAQ}, \textbf{LabVIEW}, SDR}
\skillcat{Programación}{\textbf{Python} (pandas, numpy, matplotlib), C/C++, TypeScript, FastAPI, Next.js, Docker, VBA/Excel}
\skillcat{Aeroespacial / Seguridad}{Seguimiento satelital LEO, ciberseguridad satelital, comunicaciones 6G, hábitats análogos}
\skillcat{CAD / CAE}{SolidWorks, ANSYS, FEA, simulación térmica, planos técnicos, metrología}
}

% ── PUBLICACIONES CIENTÍFICAS ─────────────────────────────────────────────────
\section{Publicaciones Científicas}

{\small\color{gristexto}
\textbf{Internacional:}\\
Sánchez González, A., Manríquez Chávez, Y., Rosas Palacios, A., \textbf{Molina Arce, J.} (2024). \textit{31st IAA Symposium on Small Satellite Missions}. IAC 2024, Milán, Italia.
\href{https://doi.org/10.52202/078365-0120}{DOI: 10.52202/078365-0120}
\vspace{2pt}
\textbf{Nacional:}\\
\textbf{Molina Arce, J.} (2024). Caracterizaci\'{o}n, dise\~{n}o y simulaci\'{o}n de un banco de pruebas mec\'{a}nicas para la medici\'{o}n de par\'{a}metros operativos. \textit{Revista Hacia el Espacio}, M\'{e}xico.
}

% ── PARTICIPACIÓN INTERNACIONAL ───────────────────────────────────────────────
\section{Participaci\'{o}n Internacional \& Reconocimientos}

{\small\color{gristexto}
\begin{itemize}
  \item \textbf{Ponente / Coautor} --- IAC 2024 (oct. 2024), Mil\'{a}n, Italia. Presentaci\'{o}n internacional en ingl\'{e}s ante audiencia cient\'{i}fica global.
  \item \textbf{Jefe de Staff} --- ICASST 2023 $\cdot$ \textbf{Comit\'{e} Organizador} --- ICASST 2022, Ciudad de M\'{e}xico.
  \item \textbf{Demostraci\'{o}n tecnol\'{o}gica} --- TV Azteca M\'{e}xico (2024): tecnolog\'{i}as Co.De Aerospace.
\end{itemize}
}

% ── IDIOMAS ───────────────────────────────────────────────────────────────────
\section{Idiomas}

{\small\color{gristexto}
\textbf{Espa\~{n}ol:} Nativo \quad \textbf{Ingl\'{e}s:} Avanzado conversacional --- apto para entrevistas t\'{e}cnicas, reuniones de trabajo y presentaciones en EE.UU.
}

% ── PIE ───────────────────────────────────────────────────────────────────────
\vspace{3pt}
{\color{lineagris}\rule{\linewidth}{0.3pt}}
{\footnotesize\centering\color{gristexto}
Disponibilidad inmediata para reubicarme en EE.UU. \textbf{$\cdot$} Sin restricciones para proceso de colocaci\'{o}n \textbf{$\cdot$} Autorizo uso de datos para fines de reclutamiento
}

\end{document}
